\documentclass[10pt,a4paper]{amsart}
\usepackage[margin=19mm]{geometry}
\usepackage{amsmath,amssymb,mathtools}
\usepackage[scr=rsfs,cal=euler]{mathalfa}
\usepackage{xfrac}
\usepackage[unicode,hidelinks,bookmarks=false]{hyperref}
\newcommand{\ie}{i.e.\ }
\newcommand{\cf}{cf.\ }
\newcommand{\resp}{respectively\ }
\newcommand{\Subsection}{\subsection}
\providecommand{\nobreakdash}{\nobreak\hbox{-}}
\setlength{\emergencystretch}{2em}
\multlinegap0pt
\usepackage{bm}
\usepackage{bbm}
\usepackage{dsfont}
\usepackage{xspace}
\usepackage{cancel}
\usepackage{mathtools}
\usepackage[shortlabels]{enumitem}
\usepackage{amsthm}
\makeatletter
\@ifundefined{prop}{\newtheorem{prop}{Proposition}}{}
\@ifundefined{lemma}{\newtheorem{lemma}[prop]{Lemma}}{}
\makeatother
\DeclareMathOperator{\dvol}{dvol}
\newcommand{\Aa}{\textbf{a}}
\newcommand{\Bb}{\textbf{b}}
\newcommand{\CC}{\mathbb{C}}
\newcommand{\Cc}{\textbf{c}}
\title[Improved fractal Weyl bounds for convex cocompact hyperbolic]{Improved fractal Weyl bounds for convex cocompact hyperbolic surfaces and large resonance-free regions}
\author{Louis Soares}
\date{Source: arXiv:2301.03023v2 (CC BY 4.0)}
\begin{document}
\maketitle

\noindent\emph{Source and licence.} Improved fractal Weyl bounds for convex cocompact hyperbolic surfaces and large resonance-free regions. Version arXiv:2301.03023v2 (CC BY 4.0), distributed under the Creative Commons Attribution 4.0 International licence, as recorded in the arXiv licence metadata for this exact version and shown on the abstract page https://arxiv.org/abs/2301.03023v2. Full rights notice: licenses/arxiv-2301.03023-CC-BY-4.0.txt.\par\smallskip

\noindent\textit{Editorial scope.} Counted unit: the Prop whose source label is \texttt{prop:main\_estimate} in the article (source0.tex lines 1809--1816, source label \texttt{prop:main\_estimate}), delivered together with the article's own complete original proof of it (source0.tex lines 1842--2028, one \texttt{proof} environment of 187 lines). No mathematics is written, repaired or re-derived: statement and proof are the article's own text line for line. The proof is detached from the statement in the source: the article states the result at lines 1809--1816 and prints its proof at lines 1842--2028, and the proof environment's own header names the statement, so the association is the article's own. Required context retained uncounted: lem:boundsderivatives (lines 825--848); part:almostMultip (lines 827--847); part:mirrorEst (lines 827--847); lem:YZ (lines 850--868); refOpRefSpace (lines 922--924); lem:contractionInH (lines 930--936); lem:welldefinedness (lines 950--966); lem:crucial\_estimate (lines 1035--1040); lem:UpperBoundBergman (lines 1138--1147); lem:separation (lines 1169--1176); phase (lines 1215--1217); Dab (lines 1224--1226); prop:FOURIER (lines 1316--1328); prop:HSnorm (lines 1664--1680); lem:weird\_sum (lines 1740--1746). Every definition, display and label of those lines is retained unchanged. Omitted from this package and not redistributed: 1--824, 848--849, 869--921, 925--929, 937--949, 967--1034, 1041--1137, 1148--1168, 1177--1214, 1218--1223, 1227--1315, 1329--1663, 1681--1739, 1747--1808, 1817--1841, 2029--2035. Nothing inside a retained excerpt is omitted or altered: every retained body is delivered exactly as the article's lines are written, so no omission receipt of any kind accompanies this package. Citations: the retained excerpts carry no citation command at all, so no bibliography entry of the article is transcribed and no reference is redistributed. Reference adaptation: none was needed and none was made; every reference inside the delivered unit resolves inside it, so no unresolved reference or citation is produced. Printed slips kept literal: no printed word, symbol, hypothesis or number of the counted statement or of its proof was altered. Layout and class: the article's own class is \texttt{amsart} with packages including ulem, enumerate, amssymb, amsthm, amsmath, mathrsfs, hyperref, geometry, braket, comment, tikz, forest, pgfplots, graphicx; the wrapper is the lane's pinned \texttt{amsart} editorial wrapper at 10pt on A4, which restates the theorem-like environments the retained text needs and adds the article's own macro definitions that the retained text needs (\texttt{\textbackslash Aa}, \texttt{\textbackslash Bb}, \texttt{\textbackslash CC}, \texttt{\textbackslash Cc}, \texttt{\textbackslash dvol}), with extra wrapper packages (bm, bbm, dsfont, xspace, cancel, mathtools, enumitem). The delivered numbering is the unit's own. Assets: no retained span contains a figure environment, a graphics-inclusion command, a PDF-page inclusion, a table or a TikZ picture; no image file is copied into this package, the package declares no asset dependency and no image file is redistributed with it. Licence: version arXiv:2301.03023v2 of this article is distributed under the Creative Commons Attribution 4.0 International licence, as recorded in the arXiv licence metadata for this exact version and shown on the abstract page https://arxiv.org/abs/2301.03023v2. A full rights notice is delivered at licenses/arxiv-2301.03023-CC-BY-4.0.txt. Characters: every retained excerpt is pure ASCII, so the delivered engine, XeLaTeX with its default Latin Modern text font, reports no missing character.
\begin{lemma}[Bounds for derivatives]\label{lem:boundsderivatives}
The following estimates hold with implied constants depending only on $\Gamma$:
\begin{enumerate}[(i)]
\item \label{part:ups} For all $b \in \mathcal{A}$ and all $z \in D_b$ with $\Aa \to b$, we have
\[
|\gamma'_{\Aa}(z)| \asymp \vert I_\Aa\vert \asymp \Upsilon_{\Aa}.
\]

\item For all $\Aa \in \mathcal{W}^\circ$, we have
\[
\Upsilon_{\Aa'} \asymp \Upsilon_{\Aa}.
\]

\item \label{part:mirrorEst} For all $\Aa \in \mathcal{W}^\circ$, we have
\[
\Upsilon_{\overline{\Aa}} \asymp \Upsilon_{\Aa}.
\]

\item \label{part:almostMultip} For all $\Aa, \Bb \in \mathcal{W}^\circ$ with $\Aa \to \Bb$, we have
\[
\Upsilon_{\Aa\Bb} \asymp \Upsilon_{\Aa} \Upsilon_{\Bb}.
\]
\end{enumerate}
\end{lemma}

\begin{lemma}[Estimates for $Z(\tau)$ and $Y(\tau)$]\label{lem:YZ}
For all $\tau > 0$, the following estimates hold with constants depending only on $\Gamma$:
\begin{enumerate}[(i)]
\item For all $\Aa \in Z(\tau)$ and all $\Aa \in Y(\tau)$, we have
\[
\Upsilon_{\Aa} \asymp \tau.
\]

\item \label{cardinalityzy} The cardinalities satisfy
\[
|Z(\tau)| \asymp |Y(\tau)| \asymp \tau^{-\delta}.
\]

\item \label{lengthsofwords} For all $\Aa \in Z(\tau)$ and all $\Aa \in Y(\tau)$, we have
\[
|\Aa| \asymp \log(\tau^{-1}).
\]
\end{enumerate}
\end{lemma}

\begin{equation}\label{refOpRefSpace}
\mathcal{L}_{Z,s}f(z) := \sum_{\substack{\Aa \in \overline{Z}' \\ \Aa \to b}} \gamma'_{\Aa}(z)^s f(\gamma_{\Aa}(z)), \quad \text{for } z \in \Omega_b(h).
\end{equation}

\begin{lemma}[Contraction in $\Omega(h)$]\label{lem:contractionInH}
There exist $N_0 \in \mathbb{N}$ and $h_0 > 0$, depending only on $\Gamma$, such that for all $h \in (0, h_0)$, all $b \in \mathcal{A}$, all $z \in \Omega_b(h)$, and all $\Aa \in \mathcal{W}_{\ge N_0}$ with $\Aa \to b$, we have
\[
\mathrm{dist}(\gamma_{\Aa}(z), \partial \Omega(h)) > h/2.
\]
In other words, $\gamma_{\Aa}(\Omega_b(h)) \subset \Omega(h/2)$.
\end{lemma}

\begin{lemma}[Well-defined action on refined space]\label{lem:welldefinedness}
Let $N_0 \in \mathbb{N}$ and $h_0 > 0$ be as in Lemma~\ref{lem:contractionInH}. For every finite set $Z \subset \mathcal{W}_{\ge N_0}$ and every $h \in (0, h_0)$, the operator
\begin{equation}\label{defi:refTOonRefS}
\mathcal{L}_{Z,s} \colon H^2(\Omega(h)) \to H^2(\Omega(h)), \quad s \in \CC
\end{equation}
defined by Equation~\eqref{refOpRefSpace}, is well-defined and of trace class.

Moreover, there exists $\tau_0 > 0$ such that for all $\tau \in (0, \tau_0)$ and all $h \in (0, h_0)$, the estimate
\[
\mathrm{dist}(\gamma_{\Aa}(z), \partial \Omega(h)) > h/2
\]
holds for all $b \in \mathcal{A}$ and all $\Aa \in Y(\tau)$ with $\Aa \to b$. Consequently, the $\tau$-refined operator
\begin{equation}\label{defi:refTOonRefS_2}
\mathcal{L}_{\tau,s} \colon H^2(\Omega(h)) \to H^2(\Omega(h)), \quad s \in \CC
\end{equation}
is also well-defined and trace class.
\end{lemma}

\begin{lemma}[Bound for complex powers of derivatives]\label{lem:crucial_estimate}
There exists $C > 0$, depending only on $\Gamma$, such that for all $b \in \mathcal{A}$, all $\Aa \in \mathcal{W}^\circ$ with $\Aa \to b$, and all $s = \sigma + i t \in \CC$ with $\sigma > 0$,
\[
z \in \Omega_b(h) \Longrightarrow |\gamma_{\Aa}'(z)^s| \le (C \Upsilon_{\Aa})^{\sigma} \exp(C h |t|).
\]
\end{lemma}

\begin{lemma}[Upper bound for the Bergman kernel]\label{lem:UpperBoundBergman}
Define
\[
\mathrm{dist}(z, \partial\Omega) := \inf_{z' \in \partial\Omega} |z - z'|.
\]
Then for all $z, w \in \Omega$,
\begin{equation}\label{UpperBoundBergmanineq}
|B_\Omega(z,w)| \le \frac{1}{\pi \, \mathrm{dist}(z, \partial\Omega) \, \mathrm{dist}(w, \partial\Omega)}.
\end{equation}
\end{lemma}

\begin{lemma}[Separation Lemma]\label{lem:separation}
There exist constants $\widetilde{C} > 0$ and $h_0 > 0$, depending only on $\Gamma$, such that for all $0 < h < h_0$ and $\tau \ge \widetilde{C} h$, the following holds:

For all $b \in \mathcal{A}$, all words $\Aa, \Bb \in Y(\tau)$ with $\Aa, \Bb \to b$, and all points $z_1, z_2 \in D_b$, we have:
\[
\gamma_{\Aa}(z_1),\, \gamma_{\Bb}(z_2) \text{ lie in the same connected component of } \Omega(h) \quad \Longrightarrow \quad \Aa = \Bb.
\]
\end{lemma}

\begin{equation}\label{phase}
\Phi_{\Aa, \Bb}(z) := \overline{\mathbb{L}(\gamma_{\Bb}'(z))} -\mathbb{L}(\gamma_{\Aa}'(z)),
\end{equation}

\begin{equation}\label{Dab}
\mathcal{D}_{\Aa, \Bb} := \left( \frac{\Upsilon_{\Aa} \Upsilon_{\Bb}}{\Upsilon_{\mathsf{red}(\Aa \overline{\Bb})}} \right)^{1/2}.
\end{equation}

\begin{prop}[Key bound for averaged oscillatory integrals]\label{prop:FOURIER} Let $\varphi_{T,H}$ be as above. There exists $T_0 = T_0(\Gamma) > 0$ such that for all $T \ge T_0$, $\eta > 0$, $T^\eta \le H \le T$, $b \in \mathcal{A}$, and $\Aa, \Bb \in \mathcal{W}^\circ$ with $\Aa, \Bb \to b$, all measurable $f \colon D \to \CC$, and all $\epsilon > 0$, the following holds with $h := \frac{1}{T}$:
\[
\left\vert \int_{\Omega_b(h)} \widehat{\varphi_{T,H}}\left(\Phi_{\Aa,\Bb}(z)\right) f(z)\, \dvol(z) \right\vert \ll_{\epsilon,\eta,\Gamma} 
\begin{cases}
h^{-\delta + 2} \Vert f \Vert_{\infty, \Omega_b(h)} & \text{if } \Aa = \Bb, \\[0.5em]
h^{-\delta + 2} \mathcal{D}_{\Aa, \Bb}^{-\delta} H^{-\delta + \epsilon} \Vert f \Vert_{\infty, \Omega_b(h)} & \text{if } \Aa \neq \Bb.
\end{cases}
\]
Here, $\mathcal{D}_{\Aa, \Bb}$ is defined as in \eqref{Dab}, and
\[
\Vert f \Vert_{\infty, \Omega_b(h)} := \sup_{z \in \Omega_b(h)} \vert f(z) \vert.
\]
\end{prop}

\begin{prop}[HS-norm of $\mathcal{L}_{\tau_0, \tau_1, s}$]\label{prop:HSnorm}
Let $h,\tau_0, \tau_1 > 0$ be sufficiently small in terms of $\Gamma$. The Hilbert--Schmidt norm of the operator
\[
\mathcal{L}_{\tau_0, \tau_1, s} = \mathcal{L}_{\tau_0, s} \mathcal{L}_{\tau_1, s}\colon H^2(\Omega(h)) \to H^2(\Omega(h))
\]
is given by
\[
\Vert \mathcal{L}_{\tau_0, \tau_1, s} \Vert_{\mathrm{HS},h}^2 = \sum_{b\in \mathcal{A}}  
\sum_{\substack{  
\Aa = \Aa_0\Aa_1 \in Y(\tau_0) \times Y(\tau_1)\\  
\Bb = \Bb_0\Bb_1 \in Y(\tau_0) \times Y(\tau_1) \\  
\Aa_0 \to \Aa_1 \to b,\ \Bb_0 \to \Bb_1 \to b } }  
\int_{\Omega_b(h)} \gamma'_{\Aa}(z)^s \overline{\gamma'_{\Bb}(z)^s}  
B_{\Omega(h)}(\gamma_{\Aa}(z), \gamma_{\Bb}(z)) \dvol(z).
\]
Here, given sets $S_0, S_1 \in \mathcal{W},$ we write $\Aa = \Aa_0\Aa_1 \in S_0 \times S_1$ to mean $\Aa_i \in S_i$ for $i \in \{0,1\}$, and similarly for $\Bb$. The function $B_{\Omega(h)}(\cdot,\cdot)$ denotes the Bergman reproducing kernel of $H^2(\Omega(h))$.
\end{prop}

\begin{lemma}[Bound for $S(\alpha, \tau)$]\label{lem:weird_sum}
For all $\tau > 0$ and all $\alpha \in (0,\delta)$, we have
\begin{equation}\label{eq:weird_sum}
S(\alpha, \tau) \ll_{\epsilon, \Gamma} \tau^{-\epsilon} \left(\tau^{-\delta} +  \tau^{- 2\delta + 2\alpha} \right)
\end{equation}
for any $\epsilon > 0$.
\end{lemma}

\begin{prop}[Main Technical Estimate]\label{prop:main_estimate}
There exists $T_0 = T_0(\Gamma) > 0$ such that for all $T \ge T_0$, $\eta > 0$, $T^{\eta} \le H \le T$, $\sigma \ge 0$, $\epsilon > 0$, and $s \in \CC$ with $\mathrm{Re}(s) \ge \sigma$, $|\mathrm{Im}(s)| \le T$, the following holds: there are parameters $h, \tau_0, \tau_1$ such that
\[
\frac{1}{2H} \int_{T-H}^{T+H} \Vert \mathcal{L}_{\tau_0, \tau_1, s+it} \Vert_{\mathrm{HS}, h}^2 \, dt 
\ll_{\epsilon,\eta ,\Gamma} C^{\sigma} H^{-(2\sigma-\delta)+\epsilon} T^{2\delta - 2\sigma}.
\]
Moreover, this bound is achieved by choosing $h = T^{-1}$ and resolution parameters $\tau_0 = c_0 T^{-1}$ and $\tau_1 = c_1 H^{-1}$ with constants $c_0, c_1 > 0$ depending only on $\Gamma$.
\end{prop}

\begin{proof}[Proof of Proposition \ref{prop:main_estimate}]
Let $h, \tau_0, \tau_1 > 0$ be sufficiently small such that the operator
\begin{equation}\label{opfinal}
\mathcal{L}_{\tau_0, \tau_1, s} \colon H^2(\Omega(h)) \to H^2(\Omega(h))
\end{equation}
is well-defined. The parameters $\tau_0$ and $\tau_1$ will be chosen in terms of $T$ and $H$ during the proof.

By Proposition \ref{prop:HSnorm}, we have
$$
\Vert \mathcal{L}_{\tau_0, \tau_1, s} \Vert_{\mathrm{HS},h}^2 = \sum_{b \in \mathcal{A}} 
\sum_{\substack{
\Aa = \Aa_0 \Aa_1 \in Y(\tau_0)\times Y(\tau_1)\\
\Bb = \Bb_0 \Bb_1 \in Y(\tau_0)\times Y(\tau_1)\\
\Aa_0 \to \Aa_1 \to b,\; \Bb_0 \to \Bb_1 \to b
}}
\int_{\Omega_b(h)} g_{\Aa,\Bb}(z;s) \,\mathrm{dvol}(z),
$$
where 
$$
g_{\Aa,\Bb}(z;s) := \gamma'_\Aa(z)^s \overline{\gamma'_\Bb(z)^s} B_{\Omega(h)}(\gamma_\Aa(z), \gamma_\Bb(z)).
$$
Note that for all $t\in \mathbb{R}$ we have
$$
g_{\Aa, \Bb}(z;s+it) = g_{\Aa, \Bb}(z;s) e^{-it  \Phi_{\Aa,\Bb}(z)},
$$
where $ \Phi_{\Aa,\Bb}$ is the phase defined in \eqref{phase}, whence
$$
\Vert \mathcal{L}_{\tau_0, \tau_1, s+it} \Vert_{\mathrm{HS},h}^2 = \sum_{b \in \mathcal{A}} 
\sum_{\substack{
\Aa = \Aa_0 \Aa_1 \in Y(\tau_0)\times Y(\tau_1)\\
\Bb = \Bb_0 \Bb_1 \in Y(\tau_0)\times Y(\tau_1)\\
\Aa_0 \to \Aa_1 \to b,\; \Bb_0 \to \Bb_1 \to b
}}
\int_{\Omega_b(h)} g_{\Aa,\Bb}(z;s) e^{-it  \Phi_{\Aa,\Bb}(z)} \,\mathrm{dvol}(z),
$$
Let $\varphi \in C^\infty(\mathbb{R})$ be a non-negative bump function supported on $[-2,2]$ with $\varphi \equiv 1$ on $[-1,1]$, and define
\[
\varphi_{T,H}(t) := \frac{1}{H} \, \varphi\left( \frac{t - T}{H} \right).
\]
Then $\varphi_{T,H}$ is supported on $[T - 2H, T + 2H]$ and equals $1/H$ on $[T - H, T + H]$, so that
\[
\frac{1}{H} \int_{T - H}^{T + H} \left\| \mathcal{L}_{\tau_0, \tau_1, s + it} \right\|_{\mathrm{HS}, h}^2 \, dt
\le \int_{-\infty}^{\infty} \varphi_{T,H}(t) \left\| \mathcal{L}_{\tau_0, \tau_1, s + it} \right\|_{\mathrm{HS}, h}^2 \, dt
=:\Sigma.
\]
Expanding the HS-norm using the previous formula, we find
\begin{equation}\label{eq:sigma-expansion}
\int_{-\infty}^{\infty} \varphi_{T,H}(t) \left\| \mathcal{L}_{\tau_0, \tau_1, s + it} \right\|_{\mathrm{HS}, h}^2 \, dt= \sum_{b \in \mathcal{A}} 
\sum_{\substack{
\Aa = \Aa_0 \Aa_1 \in Y(\tau_0)\times Y(\tau_1) \\
\Bb = \Bb_0 \Bb_1 \in Y(\tau_0)\times Y(\tau_1) \\
\Aa_0 \to \Aa_1 \to b, \;\Bb_0 \to \Bb_1 \to b
}}
\int_{\Omega_b(h)} g_{\Aa,\Bb}(z;s) \, \widehat{\varphi_{T,H}}(\Phi_{\Aa,\Bb}(z)) \, \mathrm{dvol}(z) := \Sigma.
\end{equation}

Let $\widetilde{C} > 0$ be as in Lemma \ref{lem:separation}. From now on, let
$$
\tau_0 = \widetilde{C} h.
$$
Then $\Aa_0 \neq \Bb_0$ implies that $\gamma_\Aa(z)$ and $\gamma_\Bb(z)$ lie in different components of $\Omega(h)$, so
$$
B_{\Omega(h)}(\gamma_\Aa(z), \gamma_\Bb(z)) = 0,
$$
Thus, only terms with $\Aa_0 = \Bb_0$ contribute, giving
\begin{equation}\label{eq:reexpr}
\Sigma =   \sum_{b \in \mathcal{A}} 
\sum_{\substack{
\Cc \in Y(\tau_0)\\
\Aa_1,\Bb_1 \in Y(\tau_1)\\
\Cc \to \Aa_1,\Bb_1 \to b
}}
\int_{\Omega_b(h)} g_{\Aa,\Bb}(z;s) \widehat{\varphi_{T,H}}(\Phi_{\Aa,\Bb}(z)) \,\mathrm{dvol}(z).
\end{equation}
For each $b\in \mathcal{A}$, define $Q_b$ as the set corresponding to the inner sum, and decompose it into diagonal and off-diagonal parts:
\begin{align*}
Q_b^{(1)} &:= \{ (\Cc\Aa_1,\Cc\Aa_1) : \Cc \in Y(\tau_0),\, 
\Aa_1\in Y(\tau_1),\, \Cc \to \Aa_1 \to b   \}, \\
Q_b^{(2)} &:= \{ (\Cc\Aa_1,\Cc\Bb_1) : \Cc \in Y(\tau_0),\, 
\Aa_1, \Bb_1\in Y(\tau_1), \Cc \to \Aa_1, \Bb_1 \to b ,\, \Aa_1 \neq \Bb_1 \}.
\end{align*}
Decompose $\Sigma$ accordingly, writing
$$
\Sigma = \Sigma^{(1)} + \Sigma^{(2)},
$$
where
$$
\Sigma^{(\ell)} := \sum_{b \in \mathcal{A}} 
\sum_{(\Aa,\Bb) \in Q_b^{(\ell)}}
\int_{\Omega_b(h)} g_{\Aa,\Bb}(z;s) \widehat{\varphi_{T,H}}(\Phi_{\Aa,\Bb}(z)) \,\mathrm{dvol}(z), \quad \ell\in \{1,2\}.
$$

Throughout, let $T\gg 1$ and assume $h=T^{-1}$, $\eta > 0$, $T^{\eta}\le H \le  T$, $\mathrm{Re}(s)\ge \sigma>0$, and $\vert \mathrm{Im}(s) \vert\le  K$. Moreover, $C>0$ is a constant depending only on $\Gamma$ whose precise value changes from place to place. All implied constants are allowed to depend on $\Gamma.$

Lemma \ref{lem:welldefinedness} ensures that if $\tau_0 > 0$ and $\tau_1 > 0$ are sufficiently small, then for all $(\Aa,\Bb)\in Q_b$ and $z\in \Omega_b(h)$ we have
$$
\mathrm{dist}( \gamma_{\Aa}(z), \partial \Omega(h) ) \ge h/2, \quad \mathrm{dist}( \gamma_{\Bb}(z), \partial \Omega(h) )\ge h/2.
$$
Combining this with Lemma \ref{lem:UpperBoundBergman}, we get
\begin{equation}\label{est:BKpointwise}
B_{\Omega(h)}(\gamma_{\Aa}(z), \gamma_{\Bb}(z) ) \le  \frac{1}{\pi\, \mathrm{dist}( \gamma_{\Aa}(z), \partial \Omega(h) )\, \mathrm{dist}( \gamma_{\Bb}(z), \partial \Omega(h) ) } \ll h^{-2}.
\end{equation}
Lemmas \ref{lem:boundsderivatives} and \ref{lem:YZ} imply that for each $b\in \mathcal{A}$ and $(\Aa,\Bb)\in Q_b$, we have 
$$
\Upsilon_{\Aa}  \asymp \tau_0 \tau_1, \quad \Upsilon_{\Bb}  \asymp \tau_0 \tau_1.
$$
Combining this with Lemma~\ref{lem:crucial_estimate}, we have for all $\mathrm{Re}(s) > 0,$
\begin{equation}
\vert \gamma'_{\Aa}(z)^{s+it} \vert \le  (C\Upsilon_{\Aa})^{\mathrm{Re}(s)} e^{C h (K+t)} \le ( C\tau_0 \tau_1 )^{\mathrm{Re}(s)} e^{C h (K+T+H)},
\end{equation}
where we used $h(K+T+H) = O(1) $, which follows from the assumptions above. Assuming $\tau_0, \tau_1$ are small enough so that $C\tau_0 \tau_1 < 1$, we obtain further for all $\mathrm{Re}(s) \ge \sigma > 0,$
\begin{equation}\label{final:bound_gamma_a}
\vert \gamma'_{\Aa}(z)^{s+it} \vert \ll ( C\tau_0 \tau_1 )^{\sigma},
\end{equation}
The same estimate holds for $\Bb$. Combining \eqref{final:bound_gamma_a} and \eqref{est:BKpointwise}, we deduce that for all $(\Aa,\Bb)\in Q_b$ and $z\in \Omega_b(h)$:
\begin{equation}\label{est:pointwisegab}
\vert g_{\Aa, \Bb}(z;s+it) \vert \ll  ( C\tau_0 \tau_1 )^{2\sigma} h^{-2}.
\end{equation}
Meanwhile, by Lemma \ref{lem:YZ} we have
\begin{equation}\label{final:bound_ytau}
\vert Y(\tau)\vert \asymp \tau^{-\delta}.
\end{equation}
Recall from our choices of $\tau_0$ and $h$ that
\begin{equation}\label{final:bound_tau0}
\tau_0 \asymp h = T^{-1}.
\end{equation}
From the diagonal case of Proposition \ref{prop:FOURIER}, we can now estimate:
\begin{align}
\Sigma^{(1)} &\ll \sum_{b\in \mathcal{A}}  
\sum_{\substack{ \Cc \Aa_1 \in Y(\tau_0)\times Y(\tau_1)  \\ 
\Cc \to \Aa_1 \to b
}} h^{-\delta + 2} \left\Vert g_{\Cc \Aa_1, \Cc \Aa_1}(\cdot;s) \right\Vert_{\infty,\Omega_b(h)}   \\
&\ll \vert Y(\tau_0)\vert \vert Y(\tau_1)\vert (C \tau_0\tau_1)^{2\sigma} h^{-\delta} &&\text{by \eqref{est:pointwisegab}}\\
&\ll ( C\tau_0 \tau_1 )^{-\delta + 2\sigma} h^{-\delta}&&\text{by \eqref{final:bound_ytau}} \\
&\ll C^\sigma T^{2\delta-2 \sigma} \tau_1 ^{-\delta + 2\sigma} &&\text{by \eqref{final:bound_tau0}} \label{final1}.
\end{align}
To estimate $\Sigma^{(2)}$, recall that $Q_b^{(2)}$ consists of pairs $(\Aa,\Bb)\in \mathcal{W}\times \mathcal{W}$ with $\Aa=\Cc\Aa_1$ and $\Bb=\Cc\Bb_1$ such that
\begin{itemize}
\item $\Cc\in Y(\tau_0)$ and $\Aa_1, \Bb_1 \in Y(\tau_1)$,
\item $\Cc \to \Aa_1, \Bb_1 \to b$, and
\item $\Aa_1\neq \Bb_1$.
\end{itemize}
Under these conditions, we have $\Aa \overline{\Bb} = \Cc\Aa_1 \overline{\Bb}_1 \overline{\Cc}$ and
$$
\Cc \to \Aa_1 \overline{\Bb}_1 \to \overline{\Cc}.
$$
In particular, $\mathsf{red}(\Aa\overline{\Bb}) = \Cc \, \mathsf{red}( \Aa_1 \overline{\Bb}_1 )\, \overline{\Cc}$ and $\Cc \to \mathsf{red}( \Aa_1 \overline{\Bb}_1 )\to \overline{\Cc}$. Applying Parts \eqref{part:mirrorEst} and \eqref{part:almostMultip} of Lemma \ref{lem:boundsderivatives}, we then find
\begin{itemize}
\item $\Upsilon_{\mathsf{red}(\Aa\overline{\Bb})}  \asymp \Upsilon_{\Cc} \Upsilon_{\mathsf{red}(\Aa_1 \overline{\Bb}_1)} \Upsilon_{\overline{\Cc}} \asymp \Upsilon_\Cc^2 \Upsilon_{\mathsf{red}(\Aa_1 \overline{\Bb}_1)} , $
\item $\Upsilon_{\Aa} \asymp \Upsilon_\Cc  \Upsilon_{\Aa_1},$ and
\item $\Upsilon_{\Bb} \asymp \Upsilon_\Cc  \Upsilon_{\Bb_1}.$
\end{itemize}
This implies that for all $(\Aa,\Bb)\in Q_b^{(2)}$
\begin{equation}\label{final:bound_Dab}
\mathcal{D}_{\Aa,\Bb} = \left(  \frac{\Upsilon_{\Aa} \Upsilon_{\Bb} }{\Upsilon_{\mathsf{red}(\Aa\overline{\Bb})}}\right)^{1/2} \asymp \left(  \frac{  \Upsilon_{\Aa_1} \Upsilon_{\Bb_1} }{  \Upsilon_{\mathsf{red}(\Aa_1 \overline{\Bb}_1)}  }\right)^{1/2} \asymp \tau_1 \Upsilon_{\mathsf{red}(\Aa_1 \overline{\Bb}_1)}^{-1/2}.
\end{equation}
Using Proposition \ref{prop:FOURIER}, we can now estimate the off-diagonal contribution as follows:
\begin{align*}
\Sigma^{(2)} &= \sum_{b\in \mathcal{A}}  
\sum_{ (\Aa,\Bb) \in Q_b^{(2)} } 
\int_{\Omega_b(h)} g_{\Aa,\Bb}(z;s) \widehat{\varphi_{T,H}}(\Phi_{\Aa,\Bb}(z)) \,\mathrm{dvol}(z)\\
&\ll_{\epsilon, \eta} \sum_{b\in \mathcal{A}}  \sum_{ (\Aa,\Bb) \in Q_b^{(2)} }  h^{-\delta + 2}  \mathcal{D}_{\Aa,\Bb}^{-\delta} H^{-\delta+\epsilon} \left\Vert g_{\Aa, \Bb}(\cdot;s) \right\Vert_{\infty,\Omega_b(h)} &&\text{by Prop. \ref{prop:FOURIER}}\\
&\ll \sum_{b\in \mathcal{A}} \sum_{ (\Aa,\Bb) \in Q_b^{(2)} } (C \tau_0 \tau_1)^{2\sigma} h^{-\delta}  \mathcal{D}_{\Aa,\Bb}^{-\delta} H^{-\delta+\epsilon} &&\text{by \eqref{est:pointwisegab}}\\
&\ll \sum_{b\in \mathcal{A}} \sum_{ (\Aa,\Bb) \in Q_b^{(2)} } (C \tau_0 \tau_1)^{2\sigma} \tau_1^{-\delta} h^{-\delta} H^{-\delta+\epsilon} \Upsilon_{\mathsf{red}(\Aa_1 \overline{\Bb}_1)}^{\delta/2} &&\text{by \eqref{final:bound_Dab}}\\
&\ll \vert Y(\tau_0)\vert (C \tau_0 \tau_1)^{2\sigma} \tau_1^{-\delta} h^{-\delta} H^{-\delta+\epsilon} \left(  \sum_{\Aa_1,\Bb_1\in Y(\tau_1)}   \Upsilon_{\mathsf{red}(\Aa_1 \overline{\Bb}_1)}^{\delta/2} \right)\\
&\ll (C \tau_0 \tau_1)^{-\delta + 2\sigma} h^{-\delta} H^{-\delta+\epsilon} \left(  \sum_{\Aa_1,\Bb_1\in Y(\tau_1)}   \Upsilon_{\mathsf{red}(\Aa_1 \overline{\Bb}_1)}^{\delta/2} \right) &&\text{by \eqref{final:bound_ytau}}\\
&\ll C^{\sigma}  T^{2\delta-2\sigma} \tau_1^{-\delta +2\sigma}   H^{-\delta+\epsilon} \left(  \sum_{\Aa_1,\Bb_1\in Y(\tau_1)}   \Upsilon_{\mathsf{red}(\Aa_1 \overline{\Bb}_1)}^{\delta/2} \right) &&\text{by \eqref{final:bound_tau0}}.
\end{align*}
The remaining sum can be estimated using Lemma \ref{lem:weird_sum} (applied with $\alpha=\delta/2$), which gives
\begin{equation}\label{final2}
\Sigma^{(2)} \ll_{\epsilon, \eta} C^{\sigma}  T^{2\delta-2\sigma}  H^{-\delta+\epsilon}  \tau_1^{-2\delta +2\sigma} 
\end{equation}
Combining the bounds \eqref{final1} and \eqref{final2}, we obtain
\begin{align*}
\Sigma &\ll_\epsilon C^\sigma T^{2\delta-2 \sigma} \tau_1 ^{-\delta + 2\sigma} +    C^{\sigma}  T^{2\delta-2\sigma}  H^{-\delta+\epsilon}  \tau_1^{-2\delta +2\sigma}   \\
&= C^\sigma T^{2\delta-2 \sigma} \tau_1^{-\delta + 2\sigma} \left( 1 + H^{-\delta+\epsilon}  \tau_1^{-\delta} \right),
\end{align*}
It remains to choose $\tau_1$ optimally. This may be done by taking 
$$
\tau_1 = c H^{-1}
$$
for some constant $c>0$ sufficiently small to ensure that the parameter $\tau_1$ is admissible in all the above estimates. Inserting this choice into the previous estimate yields
$$
\frac{1}{2H} \int_{T-H}^{T+H} \Vert \mathcal{L}_{\tau_0, \tau_1, s+it} \Vert_{\mathrm{HS},h}^2 dt \ll_{\epsilon, \eta} C^\sigma T^{2\delta-2 \sigma} H^{-(2\sigma-\delta)+\epsilon},
$$
completing the proof.
\end{proof}

\end{document}
